% ============================================================
% PHẦN 4 - CẤU HÌNH HUẤN LUYỆN
% ============================================================

\section{Cấu hình huấn luyện}
\label{sec:training_config}

Để đảm bảo \textbf{so sánh công bằng} giữa 5 mô hình, nhóm thiết lập một cấu hình huấn luyện chung thống nhất. Tất cả mô hình đều được huấn luyện trên cùng một tập dữ liệu đã chia (Mục~\ref{sec:data}), cùng pipeline tiền xử lý và augmentation, với các siêu tham số chung sau:

\subsection{Thiết lập chung}
Bảng \ref{tab:common_config} thể hiện thiết lập huấn luyện chung cho tất cả 5 mô hình
\begin{table}[H]
    \caption{Thiết lập huấn luyện chung cho tất cả 5 mô hình}
    \centering
    \begin{tabularx}{0.85\textwidth}{lX}
        \toprule
        \textbf{Tham số}        & \textbf{Giá trị}                              \\
        \midrule
        GPU                     & NVIDIA Tesla T4 (Kaggle Notebook)             \\
        Số epoch tối đa         & 10                                            \\
        Effective batch size    & 8 (một số mô hình dùng gradient accumulation) \\
        Thuật toán tối ưu       & AdamW (có weight decay)                       \\
        Early stopping patience & 5 epoch                                       \\
        Seed                    & 3407                                          \\
        Phân chia dữ liệu       & 70/15/15 (Train/Val/Test), stratified split   \\
        Tiền xử lý              & Pipeline chung (Mục~\ref{sec:data})           \\
        Augmentation            & Pipeline chung (Bảng~\ref{tab:augmentation})  \\
        \bottomrule
    \end{tabularx}
    \label{tab:common_config}
\end{table}

Mỗi mô hình có \textbf{learning rate} và \textbf{kích thước ảnh đầu vào} riêng, được chọn phù hợp với yêu cầu kiến trúc của từng mô hình. Chi tiết được trình bày trong Bảng~\ref{tab:model_specific_config}.

\subsection{Tham số riêng theo mô hình}

\begin{table}[H]
    \caption{Tham số riêng của từng mô hình}
    \centering
    \begin{tabular}{llll}
        \toprule
        \textbf{Mô hình} & \textbf{Learning rate} & \textbf{Image size} & \textbf{Ghi chú}                           \\
        \midrule
        Mask R-CNN       & $1 \times 10^{-4}$     & $256 \times 256$    & Batch size 2, grad accum 4                 \\
        YOLO26m-seg      & $8 \times 10^{-4}$     & $640 \times 640$    & Batch size 8                               \\
        RF-DETR          & $1 \times 10^{-4}$     & Mặc định thư viện   & Batch size 2, grad accum 4                 \\
        MobileSAM        & $1 \times 10^{-5}$     & $1024 \times 1024$  & Batch size 2, grad accum 4; freeze encoder \\
        Mask2Former      & $5 \times 10^{-5}$     & $384$               & Batch size 4, grad accum 2                 \\
        \bottomrule
    \end{tabular}
    \label{tab:model_specific_config}
\end{table}

\textbf{Lưu ý:}
\begin{itemize}
    \item Các mô hình có batch size nhỏ hơn 8 sẽ sử dụng \textbf{gradient accumulation} để đạt effective batch size $= 8$, đảm bảo sự công bằng trong quá trình tối ưu. Việc sử dụng gradient accumulation sẽ tránh bị lỗi GPU Out of Memory.
    \item MobileSAM đóng băng hoàn toàn Image Encoder (TinyViT) và Prompt Encoder, chỉ fine-tune Mask Decoder ($\approx 3,87$M tham số) để tận dụng khả năng trích xuất đặc trưng tổng quát đã được tiền huấn luyện.
    \item Tất cả mô hình đều sử dụng checkpoint tốt nhất trên tập Validation (theo \texttt{mAP@50:95}) để đánh giá trên tập Test.
\end{itemize}
